\clearpage
\hypertarget{native-481-january4}{}
\begin{center}
\rule{0.55\linewidth}{0.4pt}\\[2pt]
Astronomische Nachrichten. Vol. XXVI. No. 623. p. 367. 17 January 1848.\\[2pt]
\rule{0.55\linewidth}{0.4pt}
\end{center}
\begin{center}[\textsc{Gauss} to \textsc{Schumacher}.]\quad Göttingen, 4 January 1848.\end{center}

After a long interruption, the sky cleared yesterday, though only briefly. This allowed me to make another meridian observation of \emph{Flora}, which I give you here.

\begin{center}
Observation of Flora.\\[3pt]
\setlength{\tabcolsep}{4pt}
\begin{tabular}{l|r|r|r}
1848 & Göttingen mean time & Right ascension & Declination \\\hline
January 3 & $9^{\rm u}\ 20^{\rm m}\ 4^{\rm s}2$ & $62^\circ\ 37'\ 13''8$ & $+15^\circ\ 59'\ 24''6$ \\
\end{tabular}
\end{center}

Professor \textsc{Goldschmidt}’s elements differ from this (according to his calculation) by

\begin{center}
$+18''2$ in right ascension.\qquad $+2''4$ in declination.
\end{center}

